% ---------------------------------------------------------------------------
%  Preambolo comune alle due versioni della nota.
%  Il file principale carica babel con la lingua e definisce \nomeprop,
%  \nomelemma e \nomeoss prima di includere questo file.
% ---------------------------------------------------------------------------
\usepackage[T1]{fontenc}
\usepackage[utf8]{inputenc}

\usepackage{lmodern}
\usepackage{textcomp}
\usepackage{microtype}

\usepackage[a4paper,top=2.7cm,bottom=2.9cm,left=2.9cm,right=2.9cm]{geometry}

\usepackage{amsmath}
\usepackage{amssymb}
\usepackage{amsthm}
\usepackage{mathtools}
\usepackage{graphicx}
\usepackage{xcolor}
\usepackage{booktabs}
\usepackage{array}
\usepackage{enumitem}
\usepackage{caption}
\usepackage{titlesec}
\usepackage[hidelinks]{hyperref}

% --- tavolozza, la stessa dell'edizione integrale ----------------------------
\definecolor{inchiostro}{HTML}{1B1F23}
\definecolor{blunotte}{HTML}{1F4E79}
\definecolor{rossoterra}{HTML}{B3352B}
\definecolor{verdesalvia}{HTML}{2A7F62}
\definecolor{grigiocarta}{HTML}{F4F5F7}
\definecolor{grigiotratto}{HTML}{C8CCD2}
\definecolor{grigiotesto}{HTML}{6B7280}

\color{inchiostro}

% --- titolazione ------------------------------------------------------------
\titleformat{\section}
  {\normalfont\sffamily\Large\bfseries\color{blunotte}}{\thesection}{0.7em}{}
\titleformat{\subsection}
  {\normalfont\sffamily\large\bfseries\color{inchiostro}}
  {\thesubsection}{0.6em}{}
\titleformat{\paragraph}[runin]
  {\normalfont\sffamily\bfseries\color{inchiostro}}{}{0pt}{}

\captionsetup{font=small,labelfont={sf,bf,color=blunotte},
              labelsep=period,width=0.92\linewidth}

\graphicspath{{../figures/}}

% --- enunciati ----------------------------------------------------------------
\theoremstyle{plain}
\newtheorem{prop}{\nomeprop}[section]
\newtheorem{lemma}[prop]{\nomelemma}
\theoremstyle{remark}
\newtheorem*{oss}{\nomeoss}

% --- riferimenti alle righe dei quaderni --------------------------------------
% \mbox: un riferimento non va mai spezzato a fine riga
\newcommand{\riga}[2]{\mbox{\texttt{#1}\,\#{}#2}}
\newcommand{\righe}[3]{\mbox{\texttt{#1}\,\#{}#2--#3}}

% --- notazione ----------------------------------------------------------------
\newcommand{\dd}{\mathrm{d}}
\newcommand{\ee}{\mathrm{e}}
\newcommand{\ii}{\mathrm{i}}
\DeclareMathOperator{\sgn}{sgn}
\DeclareMathOperator{\artanh}{artanh}
\DeclareMathOperator{\gd}{gd}
\DeclareMathOperator{\sech}{sech}
\DeclareMathOperator{\sinc}{sinc}
\DeclareMathOperator{\Real}{Re}
\DeclareMathOperator{\Imag}{Im}

\setlength{\emergencystretch}{2em}
